\chapter{Structural Scale Reconstruction and Multi-Block VTK Export}

\section{Motivation}

The continuum exporter described in earlier chapters is not sufficient for
structural finite elements.  Beams and shells do not store a full
three-dimensional stress or strain field at every spatial point.  Instead,
their formulation \emph{degenerates} the continuum kinematics into a smaller
set of generalized measures:
\begin{itemize}
  \item beam sections carry axial strain, curvature, shear strains, and
        torsion rate;
  \item shell sections carry membrane strains, curvatures, and transverse
        shears.
\end{itemize}

For the present work, reconstructing the lost three-dimensional field is not a
mere visualisation problem.  The same machinery will later be used for
\emph{scale reconstruction}: transferring a reduced structural solution into an
initial condition for a continuum discretisation.  For that reason, the
reconstruction logic must live in its own compile-time module rather than be
buried inside a VTK writer.

\section{Architectural Split}

The implementation introduced in this phase is organised in three layers:
\begin{enumerate}
  \item \textbf{Constitutive snapshot layer}.  The type-erased
        \texttt{Material<Policy>} now exports a lightweight
        \texttt{SectionConstitutiveSnapshot}.  This is the structural analogue
        of the previously introduced \texttt{InternalFieldSnapshot}.  It
        carries only the constitutive data needed to reconstruct a section:
        beam stiffness parameters $(E,G,A,I_y,I_z,J,k_y,k_z)$, shell thickness
        parameters $(E,\nu,G,t,\kappa)$, and, when applicable, the current
        fiber field.
  \item \textbf{Reduction / reconstruction layer}.  The new
        \texttt{StructuralReductionPolicy<ElementT>} specialisations in
        \texttt{src/reconstruction/StructuralFieldReconstruction.hh} encode the
        formulation-specific mapping from reduced kinematics to reconstructed
        fields.  These policies are specialised for:
        \begin{itemize}
          \item \texttt{BeamElement<TimoshenkoBeam3D,3>},
          \item \texttt{TimoshenkoBeamN<N>},
          \item \texttt{ShellElement<MindlinReissnerShell3D>}.
        \end{itemize}
  \item \textbf{Carrier / export layer}.  The new
        \texttt{StructuralVTMExporter} builds a \texttt{.vtm} multi-block file
        from the reconstruction layer.  It does not implement beam or shell
        theory; it only requests samples from the reduction policy and maps
        them onto VTK carriers.
\end{enumerate}

This split is deliberate.  The exporter is allowed to disappear or to be
replaced by another post-processing backend; the reconstruction module should
remain valid because it expresses the structural theory itself.

\section{Spatial Roles and Ontology}

The reconstruction work reinforced the ontology introduced during the domain
refactor:
\begin{itemize}
  \item \texttt{Point} remains pure geometry.
  \item \texttt{Vertex} remains a geometric mesh entity.
  \item \texttt{Node} remains a carrier of degrees of freedom.
  \item \texttt{IntegrationPoint<dim>} is a spatial sampling site in the
        \emph{embedding} dimension, not in the topological dimension.
  \item \texttt{MaterialPoint} is a continuum constitutive role layered on an
        \texttt{IntegrationPoint}.
  \item \texttt{MaterialSection} is the structural analogue of
        \texttt{MaterialPoint}: a constitutive role layered on an integration
        site, but carrying generalized beam or shell measures.
  \item \texttt{NodeSection} is a nodal geometric station, not a constitutive
        point.
\end{itemize}

The crucial correction in this phase is that \texttt{MaterialSection} now binds
to \texttt{IntegrationPoint<MaterialPolicy::dim>}, i.e.\ to a point in the
embedding space.  The previous 1D-only design was semantically too narrow even
for beams, because the constitutive site of a beam is spatially embedded in
$\mathbb{R}^3$.

\section{Constitutive Snapshots}

The new \texttt{SectionConstitutiveSnapshot} is returned by a single extra
virtual method on the type-erased material wrapper.  This follows the same
reasoning as the internal-field export: the runtime indirection is paid only
when post-processing asks for it, whereas the reconstruction kernels still
operate with fully static element and field types.

\subsection{Beam Snapshot}

For elastic Timoshenko sections the snapshot contains:
\[
\{E,\; G,\; A,\; I_y,\; I_z,\; J,\; k_y,\; k_z\}.
\]
This is sufficient to reconstruct the standard sectional field:
\[
\varepsilon_{xx}(y,z) = \varepsilon_0 - z\,\kappa_y + y\,\kappa_z,
\]
\[
\sigma_{xx}(y,z) = E\,\varepsilon_{xx}(y,z),
\]
\[
\tau_{xy} = G\,\gamma_y, \qquad \tau_{xz} = G\,\gamma_z.
\]

For higher-fidelity torsion reconstruction one would need a section warping
model.  The present implementation does \emph{not} invent one.  Torsion
remains available as a generalized measure, and can later be lifted by a
specialised profile/warping policy.

\subsection{Shell Snapshot}

For Mindlin--Reissner shells the snapshot contains:
\[
\{E,\; \nu,\; G,\; t,\; \kappa\}.
\]
Given the mid-surface generalized strains
\[
\{\varepsilon_{11}, \varepsilon_{22}, \gamma_{12},
  \kappa_{11}, \kappa_{22}, \kappa_{12},
  \gamma_{13}, \gamma_{23}\},
\]
the through-thickness reconstruction is:
\[
\varepsilon_{11}(z) = \varepsilon_{11}^0 + z\,\kappa_{11}, \qquad
\varepsilon_{22}(z) = \varepsilon_{22}^0 + z\,\kappa_{22},
\]
\[
\gamma_{12}(z) = \gamma_{12}^0 + z\,\kappa_{12}, \qquad
\gamma_{13}(z) = \gamma_{13}, \qquad
\gamma_{23}(z) = \gamma_{23}.
\]
For the present linear elastic shell model, the in-plane stresses are obtained
with the plane-stress relation:
\[
\sigma_{11} = \frac{E}{1-\nu^2}\left(\varepsilon_{11} + \nu\varepsilon_{22}\right),
\qquad
\sigma_{22} = \frac{E}{1-\nu^2}\left(\nu\varepsilon_{11} + \varepsilon_{22}\right),
\]
\[
\tau_{12} = G\,\gamma_{12}, \qquad
\tau_{13} = G\,\gamma_{13}, \qquad
\tau_{23} = G\,\gamma_{23}.
\]

\subsection{Fiber Snapshot}

The previous architecture allowed nonlinear fiber sections in the solver, but
the layout and the current fiber state were hidden behind type erasure.  This
phase adds a structural export path that exposes the \emph{current} fiber field
through the snapshot:
\[
\{y_i,\; z_i,\; A_i,\; \varepsilon_i,\; \sigma_i\}_{i=1}^{n_f}.
\]
This is obtained from the current generalized strain stored at the section and
the constitutive response of each fiber material:
\[
\varepsilon_i = \varepsilon_0 - z_i\kappa_y + y_i\kappa_z,
\qquad
\sigma_i = \sigma_i(\varepsilon_i).
\]
No ad-hoc extrapolation is introduced: the fiber block reflects the actual
constitutive state of the section.

\section{Reduction Policies}

\subsection{Why Policies and Not VTK-Specific Code}

Each structural formulation degenerates the continuum in a different way.  A
Timoshenko beam, an Euler--Bernoulli beam, an MITC shell, and a future
geometrically exact beam do not lose the same kinematic information, and
therefore they do not admit the same reconstruction.  A generic ``visual
beam'' abstraction would hide this distinction and eventually become
architecturally wrong.

The chosen design is therefore:
\[
\texttt{StructuralReductionPolicy<ElementT>}
\]
with one specialisation per formulation.  This respects the origin of the
abstraction: the reconstruction belongs to the formulation, not to the
renderer.

\subsection{3D Timoshenko Beam}

For \texttt{BeamElement<TimoshenkoBeam3D,3>} and \texttt{TimoshenkoBeamN<N>},
the reduction policy reconstructs:
\begin{itemize}
  \item the axis displacement,
  \item the section rotation vector,
  \item the generalized section strain,
  \item the generalized section resultants,
  \item the reconstructed section field at any local section coordinate
        $(y,z)$.
\end{itemize}

The displacement of a reconstructed section point is computed with the standard
small-rotation lifting:
\[
\mathbf{u}(y,z) = \mathbf{u}_c + \boldsymbol{\theta}\times(0,y,z).
\]
This is the same kinematic relation that will later be needed for continuum
initialisation from a beam model.

\subsection{MITC4 Mindlin Shell}

For \texttt{ShellElement<MindlinReissnerShell3D>} the reduction policy
reconstructs:
\begin{itemize}
  \item the mid-surface displacement,
  \item the local rotation vector,
  \item the generalized shell strain at an arbitrary parametric point,
  \item the generalized shell resultants at a material site,
  \item the through-thickness field at an arbitrary signed thickness position.
\end{itemize}

The displaced position of a point offset by $z$ from the mid-surface is:
\[
\mathbf{u}(z) = \mathbf{u}_m + \boldsymbol{\theta}\times(0,0,z).
\]
Again, this is not a visual trick; it is the actual lifting implied by the
Mindlin kinematics used by the element.

\section{Compile-Time and HPC Considerations}

The new module intentionally favours compile-time structure over generic
runtime indirection:
\begin{itemize}
  \item the reconstruction layer is specialised by concrete element type;
  \item element-local vectors remain fixed-size Eigen objects;
  \item no new virtual dispatch appears inside the hot reconstruction formulas;
  \item the only extra virtual call is the constitutive section snapshot, used
        outside the solver kernel;
  \item the element DoF index caches were made \texttt{mutable}, so the
        exporter and reconstruction layer can extract element-local states in
        \texttt{const} traversals without breaking value semantics or
        reintroducing non-const APIs.
\end{itemize}

This keeps the intended performance hierarchy intact:
\begin{enumerate}
  \item solver kernels remain governed by the concrete element type and Eigen
        stack objects;
  \item post-processing may pay a small dynamic cost at the boundary to the
        type-erased material wrapper;
  \item VTK-specific costs are isolated in the exporter only.
\end{enumerate}

\section{Section Profiles}

To separate structural theory from rendering carriers, a second layer of
abstraction is introduced: \texttt{SectionProfile}.

\subsection{Beam Profiles}

For beams, the exporter currently provides:
\begin{itemize}
  \item \texttt{RectangularSectionProfile<SamplesPerEdge>},
  \item \texttt{CircularSectionProfile<Segments>}.
\end{itemize}

These profiles do not define a constitutive law.  They only provide the local
boundary points used to build the visualization carriers.  The profile points
are generated from \texttt{consteval} reference rings and scaled at runtime by
the actual section dimensions.  This keeps the topology static and the
per-element work minimal.  The VTK exporter constrains its beam and shell
carrier types with dedicated compile-time concepts:
\texttt{BeamSectionProfileLike} and \texttt{ShellThicknessProfileLike}.  This
keeps the carrier contract explicit and prevents the exporter from turning into
an untyped runtime adapter.

\subsection{Shell Thickness Profile}

For shells, the relevant geometric carrier is the thickness direction.  The
new \texttt{ShellThicknessProfile<Samples>} provides compile-time sample
locations on the interval $[-\tfrac12,\tfrac12]$.  The exporter uses this both
for the section block and for any future scale transfer through the thickness.

\section{The Multi-Block \texttt{.vtm} Contract}

The exporter writes a ParaView/VTK multiblock file with the following semantic
blocks:
\begin{enumerate}
  \item \texttt{axis}: reduced carrier of the structural model.
        For beams this is the centerline polyline; for shells it is the
        mid-surface wireframe.
  \item \texttt{surface}: a display carrier built from the section profile
        (beams) or the top/bottom shell surfaces.
  \item \texttt{sections}: section stations.  For beams this contains section
        rings; for shells this contains through-thickness section lines.
  \item \texttt{material\_sites}: actual constitutive sites.
  \item \texttt{fibers}: actual fibers when the section model provides them.
\end{enumerate}

Every geometric carrier that is meant to deform in ParaView carries a
\texttt{displacement} vector in \texttt{PointData}.  Thus the standard
\texttt{Warp By Vector} filter works directly, without any native VTK
quadrature metadata.

\subsection{A Note on Empty Trailing Blocks}

The VTK XML multiblock writer may elide a trailing empty block when the last
dataset has no geometry.  This is a writer/reader behaviour, not a loss of
semantic information in the implementation.  The exporter therefore keeps the
logical block layout fixed, but consumers should not rely on a trailing empty
block being materialised on disk when, for example, no fiber section exists in
the current model.

\section{Testing Strategy}

Two new tests were introduced:
\begin{itemize}
  \item \texttt{test\_structural\_reconstruction.cpp} verifies the
        formulation-level reconstruction:
        \begin{itemize}
          \item a 3D Timoshenko beam reconstruction is checked against the
                expected axial-plus-bending field;
          \item a flat MITC4 shell reconstruction is checked for linear
                through-thickness variation.
        \end{itemize}
  \item \texttt{test\_structural\_vtm\_export.cpp} verifies the VTK contract:
        the axis, surface, sections, and material-sites blocks exist and carry
        the expected arrays for both a beam and a shell model, while the fiber
        block is emitted only when it has actual content.
\end{itemize}

The multiblock writer was also given an auxiliary method,
\texttt{write\_with\_local\_states()}, so that the export contract can be
tested without requiring a PETSc/MPI runtime.  This turns out to be useful
beyond testing: it decouples the structural carriers from a specific solver
backend and makes the exporter equally applicable to offline reconstructed
states.  In the same spirit, the beam surface and section carriers now use the
first and last section snapshots at the element ends instead of blindly
reusing the first material site on both sides.  This is a small but important
correction for future non-uniform beam sections.

\section{Limitations and Future Work}

This phase intentionally does not solve every structural post-processing
problem.  The most relevant open items are:
\begin{itemize}
  \item Euler--Bernoulli and geometrically exact beam policies still need their
        own reduction specialisations.
  \item Torsional warping reconstruction is not yet implemented.
  \item The beam section block currently uses end and material stations; a
        future step should expose exact interpolation nodes for arbitrary
        higher-order line elements.
  \item For nonlinear shells with layered constitutive laws, a richer shell
        section snapshot will be required.
  \item Fiber blocks currently export the actual fiber field but do not yet
        generate a finite-area carrier; they are emitted as point sites.
\end{itemize}

These are acceptable limitations because the key architectural property is now
in place: \emph{reconstruction belongs to the structural formulation and is
available independently of the visualization backend}.
